\chapter{METODOLOGIA E DESENVOLVIMENTO}

O CreditMaster foi idealizado sob uma arquitetura híbrida focada em segurança, eficiência e precisão para o setor financeiro. O diferencial técnico reside no fato da operação de maneira autônoma e 100\% local (\textit{on-premise}), assegurando a proteção de dados sigilosos e o cumprimento de marcos legais como LGPD, mitigando a dependência de APIs em nuvem.

\section{Arquitetura Híbrida do Sistema}
O núcleo de tomada de decisão do sistema opera através da composição sinérgica de duas técnicas de IA generativa:

\subsection{Ajuste Fino (QLoRA) com Llama 3.1 8B}
O modelo base selecionado foi o \textbf{Llama 3.1 8B Instruct}. A modificação estrutural de seus parâmetros foi viabilizada pela utilização da biblioteca \textbf{Unsloth}, e o modelo final foi exportado no formato \textbf{GGUF} para execução local otimizada. 
Esta etapa atingiu os seguintes propósitos:
\begin{itemize}
    \item Internalização de jargões técnicos do meio bancário e financeiro.
    \item Padronização rigorosa do formato sintático dos pareceres de análise de crédito.
    \item Otimização extrema no uso de memória para permitir inferências ultrarrápidas em \textit{hardware} local, especificamente utilizando o ecossistema CUDA Toolkit em GPUs NVIDIA.
\end{itemize}

\subsection{Motor RAG (\textit{Retrieval-Augmented Generation})}
\textbf{[ANOTAÇÃO: RAG não foi implementado ainda.]}

\section{Construção do Conjunto de Dados (\textit{Dataset})}
A qualidade do modelo \textit{fine-tuned} depende diretamente da qualidade dos dados. O repositório do CreditMaster estrutura-se sobre dados proprietários contendo:
\begin{itemize}
    \item \textbf{Políticas de Crédito:} Documentos que definem diretrizes de concessão e risco de mercado (\texttt{credit\_policy\_pt.md}). \textbf{[ANOTAÇÃO: A politica vai ser implementada ao modelo junto com o RAG.]}
    \item \textbf{Histórico de Pareceres:} Uma base sintética contendo milhares de exemplos rotulados (\texttt{pareceres\_credito\_2000.jsonl}), construída através de metodologias automatizadas de roteirização (\texttt{gerar\_pareceres.py}). \textbf{[ANOTAÇÃO: Como não tive acesso a pareceres reais, utilizei uma base de dados da Lending Club disponibilizadas via Kaggle e gerei pareceres sintéticos.]}
\end{itemize}
A base de dados foi construída com base em 14 variáveis centrais: valor do empréstimo (\texttt{loan\_amnt}), prazo (\texttt{term}), tempo de emprego (\texttt{emp\_length}), tipo de moradia (\texttt{home\_ownership}), renda anual (\texttt{annual\_inc}), finalidade (\texttt{purpose}), relação dívida/renda (\texttt{dti}), pontuação mínima de crédito (\texttt{fico\_range\_low}), atrasos nos últimos 2 anos (\texttt{delinq\_2yrs}), consultas nos últimos 6 meses (\texttt{inq\_last\_6mths}), utilização do limite rotativo (\texttt{revol\_util}), número de contas ativas abertas (\texttt{open\_acc}), ocorrência de registros públicos (\texttt{pub\_rec}) e o status histórico do contrato (\texttt{loan\_status}).

Os dados foram convertidos por meio de rotinas em Python (\texttt{gerar\_pareceres.py}) para o formato padrão de conversação \textit{Chat Template} do Llama 3.1 no formato JSONL (\texttt{pareceres\_credito\_2000.jsonl}), contendo as mensagens de sistema (\textit{system prompt}), a requisição estruturada do usuário com as métricas financeiras do cliente e a conclusão técnica gerada em formato JSON (\textit{assistant}). Do total de 1.998 amostras processadas, 1.698 (85\%) foram alocadas para o conjunto de treinamento e 300 (15\%) para o conjunto de validação e teste.

\section{Visão Geral das Tecnologias Aplicadas}
Para suportar o rigor da arquitetura, o \textit{Tech Stack} consolidou as seguintes soluções:
\begin{itemize}
    \item \textbf{Camada de IA e Modelagem:} Python 3.10+, Llama 3.1 8B, Unsloth, QLoRA e exportação GGUF.
    \item \textbf{Motor de Inferência:} Ambiente Linux Mint suportando Ollama e CUDA Toolkit.
    \item \textbf{Pipeline de Dados:} Pandas. \textbf{[ANOTAÇÃO: Aqui serão também, adicionados o utilizado para o RAG]}
\end{itemize}
